% Based on templates/lecture.tex; actual session 22, 2026-11-05 (Week 12).
% IO and Disk.pdf, slides 10.1--10.19; video titles only, no video content read.
\begingroup
\renewcommand{\recordid}[1]{}
\renewcommand{\relatedexercise}[2]{\hyperref[exercise:#1]{#2}}
\section{第 22 节课：I/O 系统与磁盘读写}
\label{lecture:SC2005-L22}

\lecturemeta
  {SC2005-L22}
  {2026-11-05（第十二教学周）}
  {IO and Disk，Part 10: I/O System \& Disk}
  {10.1--10.19（PDF pp.~1--19）；I/O 管理、磁盘结构与读写时间}
  {Intro and IO System Design and Implementation；Intro and Disk Structure；仅按标题对应，视频未读取}
  {已完成学习并要求整理；讲义与例题已核对}

\linespread{1.15}\selectfont
\setlength{\parskip}{0.4em}

\keyidea{文件系统确定要访问的磁盘块，I/O 子系统协调请求与数据传输，驱动和控制器操作设备。机械硬盘的读写时间包含寻道、旋转等待和传输；顺序读取通过减少重复定位提高效率。}

\lecturetopic{SC2005-L22-T01}{I/O 分层与设备管理}

\textbf{来源：10.1--10.6。}设备在数据单位（字符流或块）、访问方式（顺序或随机）、传输时序、独占或共享、速度与读写方向上存在差异。I/O 系统的两个设计目标是 efficiency（效率），减少设备等待与处理开销；以及 generality（通用性），向程序提供统一接口，由下层处理硬件差异。

\begin{center}
% Original schematic based on IO and Disk.pdf, slide 10.5.
\begin{tikzpicture}[
  font=\footnotesize,>=Latex,
  layer/.style={draw=noteBlue,fill=noteBlue!5,rounded corners=2pt,
    minimum width=2.2cm,minimum height=0.85cm,align=center}
]
  \node[layer] (app) at (0,0) {应用程序\\I/O 请求};
  \node[layer] (kernel) at (2.9,0) {内核 I/O 子系统\\通用管理};
  \node[layer] (driver) at (5.8,0) {驱动程序\\设备专用软件};
  \node[layer,fill=gray!10] (controller) at (8.7,0) {控制器\\硬件};
  \node[layer,fill=gray!10] (device) at (11.6,0) {设备\\实际读写};
  \draw[->] (app) -- (kernel);
  \draw[->] (kernel) -- (driver);
  \draw[->] (driver) -- (controller);
  \draw[->] (controller) -- (device);
\end{tikzpicture}


{\small 据讲义 10.5 原创重绘；箭头表示请求逐层下达，完成状态由底层向上报告。}
\end{center}

\textbf{Device driver（设备驱动程序）是软件}，将通用请求转换为设备命令，处理状态和错误；\textbf{device controller（设备控制器）是硬件}，控制实际设备执行操作。Bus（总线）连接处理器、内存与控制器。上层程序调用 \texttt{read()} 等接口，无须直接操作不同设备的控制寄存器。

\textbf{I/O scheduling（I/O 调度）}安排排队请求的服务顺序。磁盘请求的位置不同，调整顺序可以减少磁头移动。具体磁盘调度算法安排在下一节。

\lecturetopic{SC2005-L22-T02}{缓冲、缓存与假脱机}

\textbf{来源：10.7--10.8。}三种机制都可能暂存数据，区别在于用途。
\begin{center}
\begin{tabularx}{\textwidth}{@{}p{3.15cm}XX@{}}
\toprule
机制 & 主要作用 & 典型情境 \\
\midrule
Buffering（缓冲） & 协调传输双方的速度或传输单位差异 & 数据先进入内存缓冲区，再由程序取走；小份输入可组成较大的处理单位 \\
Caching（缓存） & 保留数据副本，加快后续访问 & 文件块留在内存，再次读取时减少磁盘访问 \\
Spooling（假脱机） & 保存输出作业，协调设备的逐个服务 & 多个打印任务先进入队列，打印机依次处理 \\
\bottomrule
\end{tabularx}
\end{center}

缓冲可以吸收短时间的速度差；输入长期快于消费速度时，有限缓冲区仍会装满。同一块内存可以兼具缓冲和缓存作用。打印程序提交作业后可继续执行，实际输出由 spool 队列协调。

\lecturetopic{SC2005-L22-T03}{阻塞、异步与双缓冲}

\textbf{来源：10.9--10.12。}Blocking I/O（阻塞式 I/O）在请求无法立即完成时，让调用进程进入 Waiting / Blocked 状态；I/O 完成后进入 Ready 状态，再等待 CPU 调度。\textbf{CPU 可以运行其他就绪进程}。

Asynchronous I/O（异步 I/O）在请求被接受后返回，调用者继续处理不依赖本次结果的工作，完成时再接收通知并使用结果。

\textbf{Double buffering（双缓冲）。}用两个缓冲区交替接收和交付数据：
\begin{center}
\begin{tabularx}{\textwidth}{@{}lXX@{}}
\toprule
阶段 & 设备 & 程序 \\
\midrule
启动 & 第一批读入 $b_1$ & 等待第一批数据 \\
交替一 & 第二批读入 $b_2$ & 处理 $b_1$ 中的数据 \\
交替二 & 第三批读入已空出的 $b_1$ & 处理 $b_2$ 中的数据 \\
\bottomrule
\end{tabularx}
\end{center}
读取与处理可重叠，缓冲区再次写入前必须用完旧数据。\textbf{辅助时间模型：}若每批 I/O 耗时 $T_I$，另一阶段处理耗时 $T_P$，且两阶段可独立重叠、忽略协调开销，则串行每批需 $T_I+T_P$，流水线稳定阶段的批次间隔约为 $\max(T_I,T_P)$；启动、收尾仍需时间。改善的是流水线吞吐量。

\textbf{一次读取的路径。}系统调用进入内核后，先检查请求能否直接满足（如缓存命中）；否则交给驱动，必要时阻塞调用进程。驱动向控制器下达命令，设备完成操作并产生中断，中断处理程序与驱动更新完成状态，内核交付结果并使相关等待进程就绪。系统调用、驱动执行、中断处理、必要的进程切换以及数据复制都消耗资源，故 I/O 成本不只包括设备工作时间。

\lecturetopic{SC2005-L22-T04}{磁盘结构与读写时间}

\textbf{来源：10.13--10.19；以下采用机械硬盘模型。}

\noindent\begin{minipage}[c]{0.51\textwidth}
Platter（盘片）上的同心圆为 track（磁道），每条磁道划分为 sector（扇区）；不同盘面上同半径的磁道组成 cylinder（柱面）。磁头沿半径移动以选择磁道，盘片旋转使目标扇区经过磁头。

操作系统使用一维块编号访问设备，底层将编号映射到存储位置。本讲例题取每扇区 512 B；文件系统的一个分配块可以包含多个扇区。
\end{minipage}\hfill
\begin{minipage}[c]{0.46\textwidth}
\centering
\resizebox{\linewidth}{!}{% Original top-view schematic based on IO and Disk.pdf, slides 10.15--10.16.
% An annular sector is highlighted; the head lies on the same track at angle 0.
\begin{tikzpicture}[font=\footnotesize,>=Latex]
  \fill[gray!7] (0,0) circle (1.65);
  \draw[gray] (0,0) circle (1.65);
  \draw[gray!65] (0,0) circle (0.7);
  \draw[noteBlue] (0,0) circle (1.25);
  \filldraw[fill=noteBlue!25,draw=noteBlue]
    (35:1.1) arc (35:65:1.1) -- (65:1.4)
    arc (65:35:1.4) -- cycle;
  \fill[gray!60] (0,0) circle (0.14);
  \draw[very thick] (1.25,0) -- (2.65,0);
  \fill[noteBlue] (1.18,-0.07) rectangle (1.32,0.07);
  \node[anchor=west] at (2.75,0) {磁头及磁臂};
  \node[anchor=west] (sector) at (2.3,1.3) {目标扇区};
  \draw[->] (sector.west) -- (50:1.26);
  \node[anchor=east] (track) at (-2,0.65) {磁道};
  \draw[->] (track.east) -- (150:1.25);
  \draw[->,noteBlue,thick] (85:1.85) arc (85:145:1.85);
  \node[anchor=east] at (-1.1,1.95) {盘片旋转};
  \draw[<->,thick] (0.5,-0.4) -- (2.1,-0.4);
  \node[anchor=west] at (2.3,-0.65) {寻道：沿半径移动};
\end{tikzpicture}
}

{\footnotesize 据讲义 10.15--10.16 原创示意。}
\end{minipage}\par

\textbf{三个时间分量。}Seek time（寻道时间）用于移动磁头到目标磁道；rotational latency（旋转等待时间）用于等待目标扇区；transfer time（传输时间）用于实际读写数据。\textbf{位于正确磁道仍可能需要旋转等待。}总时间为定位时间加传输时间，定位时间包含寻道与旋转等待。

设转一圈需 $T_{\mathrm{rev}}$，每磁道有 $S$ 个扇区。目标扇区角度均匀分布时，平均旋转等待为 $T_{\mathrm{rev}}/2$。若 $N$ 个连续扇区在同一磁道上、只需一次定位，则
\[
  \boxed{T_{\mathrm{I/O}}=
  \underbrace{T_{\mathrm{seek}}+\frac{T_{\mathrm{rev}}}{2}}_{T_{\mathrm{position}}}
  +\underbrace{\frac{N}{S}T_{\mathrm{rev}}}_{T_{\mathrm{transfer}}}}.
\]
跨磁道时按题意另计定位或切换成本。\textbf{计算统一采用 10.17 的参数：}$T_{\mathrm{seek}}=10$ ms、$T_{\mathrm{rev}}=8$ ms、平均等待 4 ms、$S=1024$，故一次平均定位为 14 ms。

讲义 10.18 的随机读按\textbf{每个扇区单独定位}计算，$N$ 个扇区约需 $N(14+8/1024)$ ms；10.19 的顺序读将定位成本分摊到连续数据上。KB、MB 在这些例题中采用二进制换算。完整条件和计算见\relatedexercise{SC2005-L22-E01}{分散扇区随机读取}与\relatedexercise{SC2005-L22-E02}{连续数据顺序读取}。请求大小本身不能确定定位次数，须结合数据布局判断。

\textbf{一分钟回顾。}驱动是软件，控制器是硬件，通用 I/O 子系统协调请求。缓冲解决传输不匹配，缓存加速重复访问，假脱机组织输出作业。阻塞让调用进程等待，异步与双缓冲可重叠 I/O 和处理。机械硬盘时间分为寻道、旋转等待和传输；计算时先数清定位次数，再计算连续数据的传输时间。
\endgroup
